\documentclass[11pt,a4paper]{article}
\usepackage{geometry}
\geometry{left=2cm,right=2cm,top=2.5cm,bottom=2.5cm}
\usepackage{booktabs}
\usepackage{xcolor}
\usepackage{tabularx}
\usepackage{hyperref}
\usepackage{amssymb}
\usepackage{longtable}

\hypersetup{
    colorlinks=true,
    linkcolor=blue,
    filecolor=magenta,      
    urlcolor=cyan,
}

\definecolor{primary}{HTML}{1A365D}
\definecolor{secondary}{HTML}{2B6CB0}

\usepackage{xepersian}
\settextfont{Tahoma}
\setlatintextfont{Arial}

\begin{document}

\begin{center}
    \vspace*{0.5cm}
    {\Huge\bfseries طرح پیشنهادی فنی و مالی یکپارچه‌سازی سیستم لجستیک} \\[0.4cm]
    {\Large\bfseries پیاده‌سازی وب‌سرویس‌های استعلام، تفکیک انبارها، زمان‌بندی و ثبت سفارش (ماژول H)} \\[0.8cm]
    \textbf{نسخه پیشنهادی ۱.۰} --- \textbf{منطبق بر سند تحلیل و طراحی سیستم لجستیک (۱۴۰۵ / ۲۰۲۶)} \\[1.2cm]
\end{center}

\tableofcontents
\newpage

\section{هدف و دامنه (Objective \& Scope)}
هدف این سند، تعریف و تشریح فرآیند توسعه لایه یکپارچه‌سازی فروشگاه ناپ‌کامرس با سامانه لجستیک اختصاصی است. ماژول ناپ‌کامرس به عنوان یک واسط (\lr{Adapter})، اطلاعات سبد خرید و آدرس خریدار را جمع‌آوری کرده و به سرور لجستیک ارسال می‌کند. سرور لجستیک نیز بر اساس موقعیت انبارها، موجودی، قوانین ارسال، ظرفیت پیک، تقویم و روزهای تعطیل، گزینه‌های معتبر ارسال را به صورت پویا محاسبه و برمی‌گرداند.

\subsection{خروجی مورد انتظار برای کاربر}
\begin{itemize}
    \item نمایش روش‌های ارسال مجاز (پیک اختصاصی، پست، تیپاکس، حمل ویژه باربری).
    \item نمایش تاریخ یا بازه تقریبی تحویل هر روش ارسال.
    \item نمایش و امکان انتخاب بازه‌های زمانی تحویل (مانند ۲ تیر ساعت ۱۲ تا ۱۶) در صورت لزوم.
    \item محاسبه زنده هزینه قطعی یا تخمینی ارسال.
    \item نمایش شفاف چندمرسوله‌ای بودن سفارش و لیست اقلام هر مرسوله به تفکیک انبار.
    \item نمایش پیام‌های واضح در صورت عدم امکان ارسال یک یا چند کالا به مقصد خریدار.
\end{itemize}

\subsection{موارد خارج از دامنه نسخه اول (Out of Scope)}
\begin{itemize}
    \item مسیربندی بهینه لحظه‌ای رانندگان و اپلیکیشن راننده.
    \item تسویه حساب‌های مالی با شرکت‌های حمل‌کننده خارجی.
    \item فرآیند کامل مرجوعی کالا (اگرچه زیرساخت و شناسه‌های لازم پیش‌بینی می‌شود).
\end{itemize}

\section{اصول معماری و مسئولیت سیستم‌ها}
اصل بنیادی معماری این سیستم بر پایه ایزوله‌سازی است؛ ناپ‌کامرس صرفاً نقش واسط (\lr{Adapter}) را بازی می‌کند و هیچ‌گونه منطق محاسباتی لجستیک (هزینه‌ها، پوشش شهرها یا ظرفیت بازه‌ها) در آن نگهداری نمی‌شود.

\begin{table}[htbp]
\centering
\begin{tabularx}{\textwidth}{l X}
\toprule
\textbf{سیستم} & \textbf{مسئولیت‌های اصلی} \\
\midrule
\textbf{ناپ‌کامرس (ماژول اتصال)} & خواندن سبد، مشتری و آدرس؛ تبدیل داده به قرارداد وب‌سرویس؛ نمایش گزینه‌ها؛ نگهداری موقت انتخاب کاربر؛ ثبت نهایی سفارش و دریافت وضعیت. \\
\midrule
\textbf{سامانه لجستیک} & تشخیص انبار و گروه مرسوله؛ بررسی امکان روش‌های حمل؛ محاسبه تاریخ، بازه زمانی (\lr{Time Slot}) و هزینه؛ رزرو موقت ظرفیت؛ ایجاد مرسوله؛ مدیریت وضعیت تحویل. \\
\midrule
\textbf{کاربر (خریدار)} & انتخاب آدرس، روش ارسال و بازه زمانی؛ تایید هزینه؛ ثبت سفارش و پرداخت. \\
\midrule
\textbf{مدیر فروشگاه} & مدیریت نگاشت‌ها (\lr{Mappings})، تنظیمات اتصال، مشاهده لاگ خطاها و اجرای مجدد دستی (\lr{Retry}). \\
\bottomrule
\end{tabularx}
\caption{جدول مسئولیت سیستم‌ها}
\end{table}

\noindent \textbf{قاعده تصمیم‌گیری:} فروشگاه نباید به صورت محلی تصمیم بگیرد که چه مقصدی با چه روشی و با چه هزینه‌ای ارسال شود؛ تمامی تصمیم‌ها باید از پاسخ‌های معتبر سامانه لجستیک استخراج شوند.

\section{واژگان و موجودیت‌های اصلی (Glossary)}
\begin{table}[htbp]
\centering
\begin{tabularx}{\textwidth}{l X}
\toprule
\textbf{اصطلاح} & \textbf{تعریف} \\
\midrule
\lr{Cart} & سبد خرید جاری کاربر در ناپ‌کامرس. \\
\lr{Cart Item} & یک ردیف کالا همراه با تعداد و ویژگی‌های انتخاب‌شده. \\
\lr{Warehouse} & انبار یا نقطه تأمین کالا. \\
\lr{Shipment Group} & گروهی از اقلام که قابلیت ارسال در یک مرسوله یا مسیر مشترک را دارند. \\
\lr{Delivery Option} & یک گزینه قابل انتخاب در سطح کل سبد یا یک گروه مرسوله. \\
\lr{Time Slot} & بازه تاریخ و ساعت قابل رزرو برای تحویل. \\
\lr{Quote} & نتیجه موقت استعلام ارسال با زمان انقضا. \\
\lr{Selection} & انتخاب کاربر از بین گزینه‌ها و بازه‌ها. \\
\lr{Logistics Order} & رکورد نهایی سفارش ثبت شده در سامانه لجستیک. \\
\lr{External Code} & شناسه پایدار بین دو سیستم برای شهر، استان، انبار و روش ارسال. \\
\bottomrule
\end{tabularx}
\caption{واژگان و مفاهیم اصلی لجستیک}
\end{table}

\newpage
\section{جریان کامل سبد خرید تا تحویل (End-to-End Workflow)}
\begin{enumerate}
    \item کاربر کالاها را به سبد خرید خود اضافه می‌کند.
    \item در صفحه سبد یا فرآیند تسویه حساب، آدرس تحویل انتخاب یا ثبت می‌شود.
    \item ماژول اطلاعات سبد، ویژگی فیزیکی کالاها، شناسه انبار و مقصد را آماده می‌کند.
    \item ماژول درخواست استعلام قیمت را به سامانه لجستیک ارسال می‌کند.
    \item لجستیک اقلام را بر اساس انبار مبدا و قواعد عملیاتی گروه‌بندی می‌کند.
    \item گزینه‌های معتبر شامل روش، تاریخ، بازه و هزینه بازگردانده می‌شوند.
    \item کاربر یک گزینه و در صورت نیاز بازه زمانی (\lr{Time Slot}) را انتخاب می‌کند.
    \item قبل از ثبت نهایی یا پرداخت سفارش، وضعیت نقل‌قول (\lr{Quote}) دوباره اعتبارسنجی (\lr{Validate}) می‌شود.
    \item پس از ثبت موفق سفارش، سفارش لجستیک (\lr{Logistics Order}) و مرسوله‌ها (\lr{Shipments}) ایجاد می‌شوند.
    \item تغییرات وضعیت ارسال از طریق وب‌هوک (\lr{Webhook}) یا استعلام دوره‌ای (\lr{Polling}) به فروشگاه منتقل می‌شود.
\end{enumerate}

\paragraph{نمودار متنی جریان کار (Workflow Diagram):}
\begin{latin}
\begin{verbatim}
Customer Cart + Address
        |
        v
nopCommerce Integration Plugin
        |
        | POST /delivery-quotes
        v
Logistics Engine
  - Warehouse resolution
  - Shipment grouping
  - Method eligibility
  - Capacity / calendar
  - Cost estimation
        |
        v
Delivery Options + Time Slots + Estimated Cost
        |
        v
User Selection -> Validate -> Create Logistics Order
\end{verbatim}
\end{latin}

\section{نگاشت شناسه‌ها و داده‌های مرجع (Identifier Mapping)}
برای جلوگیری از تداخل شناسه‌های داخلی دیتابیس ناپ‌کامرس با سرور لجستیک، تمامی ارتباطات بر اساس کدهای خارجی پایدار (\lr{External Codes}) نگاشت می‌شوند.

\begin{table}[htbp]
\centering
\begin{tabularx}{\textwidth}{c c c c l}
\toprule
\textbf{نوع} & \textbf{شناسه ناپ‌کامرس} & \textbf{کد خارجی پایدار} & \textbf{شناسه لجستیک} & \textbf{مثال} \\
\midrule
استان & \lr{StateProvinceId} & \lr{IR-TEH} & 8 & تهران \\
شهر & \lr{CityId} یا ویژگی آدرس & \lr{IR-TEH-TEHRAN} & 110 & تهران \\
انبار & \lr{WarehouseId} & \lr{WH-TEH-01} & 7 & انبار مرکزی تهران \\
روش ارسال & \lr{ShippingMethodId} & \lr{COURIER} & 3 & پیک اختصاصی \\
\bottomrule
\end{tabularx}
\caption{ساختار پیشنهادی نگاشت شناسه‌ها}
\end{table}

\paragraph{الزامات داده‌های مرجع لجستیک:}
\begin{itemize}
    \item وب‌سرویس همگام‌سازی استان‌ها و شهرها باید حاوی \lr{Id}، \lr{ExternalCode}، نام فارسی و وضعیت فعال باشد.
    \item موجودیت انبار باید به صورت مستقل شامل \lr{warehouseId}، \lr{externalCode} و \lr{address} بازگردد و شناسه‌ها نباید به صورت متنی در فیلدهای آدرس ادغام شوند.
    \item نگاشت‌ها باید در پنل ادمین ماژول لجستیک ناپ‌کامرس قابل مشاهده و اصلاح باشند.
    \item در صورت نبود نگاشت، سیستم باید از حدس زدن پرهیز کرده و خطای صریح بازگرداند.
\end{itemize}

\begin{latin}
\begin{verbatim}
// Sample Warehouse Address Structure in Mapping Web Service
{
  "warehouseId": 7,
  "externalCode": "WH-TEH-01",
  "name": "Tehran Central Warehouse",
  "address": {
    "provinceCode": "IR-TEH",
    "cityCode": "IR-TEH-TEHRAN",
    "postalCode": "1234567890",
    "addressLine": "Tehran, Example Street, No 4",
    "latitude": 35.70,
    "longitude": 51.40
  }
}
\end{verbatim}
\end{latin}

\section{قرارداد وب‌سرویس استعلام گزینه‌های ارسال}
\noindent \textbf{آدرس اندپوینت پیشنهادی:} \lr{POST /api/v1/integration/delivery-quotes}

\subsection{ساختار نمونه درخواست استعلام (Request)}
\begin{latin}
\begin{verbatim}
{
  "requestId": "c69492d1-63bd-4f08-91db-14a5831ef001",
  "storeId": "STORE-1",
  "cartId": "CART-9001",
  "customer": {
    "customerId": "852",
    "mobile": "09121234567"
  },
  "destination": {
    "addressId": "990",
    "receiverName": "Ali Rezaei",
    "mobile": "09121234567",
    "provinceCode": "IR-TEH",
    "provinceName": "Tehran",
    "cityCode": "IR-TEH-TEHRAN",
    "cityName": "Tehran",
    "postalCode": "1234567890",
    "addressLine": "Example Street, No 10",
    "latitude": 35.7001,
    "longitude": 51.4002
  },
  "currency": "IRR",
  "items": [
    {
      "cartItemId": "501",
      "productId": "145",
      "sku": "SKU-145",
      "productName": "Sample Product A",
      "quantity": 2,
      "warehouseCode": "WH-TEH-01",
      "unitWeightGrams": 1500,
      "dimensionsCm": { "length": 30, "width": 20, "height": 15 },
      "unitPrice": 25000000,
      "requiresShipping": true,
      "isFragile": false
    }
  ]
}
\end{verbatim}
\end{latin}

\subsection{ساختار نمونه پاسخ استعلام در سطح سبد (Response)}
\begin{latin}
\begin{verbatim}
{
  "isSuccess": true,
  "quoteId": "QT-20260721-10025",
  "expiresAtUtc": "2026-07-21T13:30:00Z",
  "currency": "IRR",
  "options": [
    {
      "optionId": "OPT-FAST-1",
      "title": "Fast Shipping",
      "totalEstimatedCost": 3200000,
      "estimatedDeliveryFrom": "2026-07-22",
      "estimatedDeliveryTo": "2026-07-22",
      "requiresTimeSlot": true,
      "timeSlots": [
        {
          "slotId": "SLOT-101",
          "date": "2026-07-22",
          "from": "12:00",
          "to": "16:00",
          "additionalCost": 0,
          "capacityStatus": "Available"
        }
      ]
    }
  ],
  "warnings": []
}
\end{verbatim}
\end{latin}

\paragraph{فیلدهای الزامی پاسخ:} \lr{quoteId} (شناسه نقل‌قول لجستیک)، \lr{expiresAtUtc} (زمان اعتبار قیمت/ظرفیت)، \lr{optionId} (شناسه گزینه ارسال)، \lr{totalEstimatedCost} (هزینه کل)، \lr{shipments} (تفکیک مرسوله‌های انبارها) و \lr{timeSlots} (بازه زمانی در صورت نیاز).

\section{مدل چندانباره و گروه‌های مرسوله}
یک سبد خرید ممکن است از چند انبار فیزیکی تامین شود. در فرانت‌اند سایت، به جای نمایش چندین انتخاب پیچیده برای هر انبار، خریدار گزینه‌های کلی سبد خرید نظیر «اقتصادی»، «سریع» و «زمان‌بندی‌شده» را مشاهده می‌کند، در حالی که در پشت صحنه جزئیات تفکیک انبارها در مرسوله‌ها نمایش داده می‌شود.

\begin{table}[htbp]
\centering
\begin{tabularx}{\textwidth}{l X}
\toprule
\textbf{سناریو} & \textbf{رفتار پیشنهادی} \\
\midrule
دو کالا از یک انبار & ارسال در یک مرسوله واحد با روش و زمان‌بندی مشترک. \\
\midrule
دو کالا از دو انبار نزدیک & در صورت امکان تجمیع کالاها (\lr{Consolidation})، گزینه‌های تک‌مرسوله‌ای و چندمرسوله‌ای پیشنهاد می‌شوند. \\
\midrule
کالای حجیم + کالای عادی & تفکیک به دو گروه مرسوله؛ ارسال کالا با حمل ویژه (باربری) و کالای عادی با پست/پیک. \\
\midrule
یک کالا غیرقابل ارسال & کل نقل‌قول ناموفق می‌شود یا پیام حذف کالا نمایش داده می‌شود تا از ثبت سفارش ناقص جلوگیری شود. \\
\bottomrule
\end{tabularx}
\caption{رفتارهای فرآیند تفکیک مرسوله‌ها}
\end{table}

\paragraph{نمونه ساختار تفکیک مرسوله‌ها (Shipment Groups Payload):}
\begin{latin}
\begin{verbatim}
{
  "shipmentGroups": [
    {
      "shipmentGroupId": "SG-1",
      "warehouseCode": "WH-TEH-01",
      "cartItemIds": ["501", "503"],
      "constraints": []
    },
    {
      "shipmentGroupId": "SG-2",
      "warehouseCode": "WH-ESF-01",
      "cartItemIds": ["502"],
      "constraints": ["FRAGILE"]
    }
  ]
}
\end{verbatim}
\end{latin}

\section{انتخاب، رزرو و اعتبارسنجی مجدد (Selections \& Validation)}
پس از انتخاب گزینه و بازه زمانی توسط کاربر، مقادیر در نشست یا دیتابیس ناپ‌کامرس ذخیره می‌شوند. هرگونه تغییر در مقادیر سبد خرید، آدرس یا مشخصات سفارش، انتخاب قبلی را منقضی کرده و استعلام جدیدی صادر می‌کند.

\paragraph{اندپوینت ثبت انتخاب موقت:} \lr{POST /api/v1/integration/delivery-selections}

\paragraph{اندپوینت اعتبارسنجی نهایی پیش از پرداخت سفارش:} \lr{POST /api/v1/integration/delivery-quotes/\{quoteId\}/validate}

\begin{latin}
\begin{verbatim}
// Sample Validation Payload
{
  "cartId": "CART-9001",
  "optionId": "OPT-FAST-1",
  "timeSlotId": "SLOT-101",
  "cartHash": "sha256:..."
}
\end{verbatim}
\end{latin}
در صورت پر شدن ظرفیت بازه یا تغییر قیمت در سیستم لجستیک، پاسخ با مقدار \lr{isValid = false} برگشته و از کاربر درخواست می‌شود گزینه‌ها را مجدداً استعلام و انتخاب کند:

\begin{latin}
\begin{verbatim}
{
  "isValid": false,
  "code": "TIME_SLOT_CAPACITY_EXHAUSTED",
  "message": "The selected time slot no longer has capacity.",
  "requiresRequote": true
}
\end{verbatim}
\end{latin}

\section{ایجاد سفارش لجستیکی و پیگیری وضعیت}
پس از اتمام موفق پرداخت، سفارش نهایی ناپ‌کامرس ثبت شده و درخواست ایجاد سفارش لجستیک با الگوی تکرارناپذیری ارسال می‌گردد:

\noindent \textbf{آدرس اندپوینت:} \lr{POST /api/v1/integration/logistics-orders} \\
\noindent \textbf{هدر الزامی:} \lr{Idempotency-Key: nop-order-78512}

\begin{latin}
\begin{verbatim}
// Request Payload for Logistics Order
{
  "storeOrderId": "78512",
  "storeOrderNumber": "ORD-20260721-78512",
  "quoteId": "QT-20260721-10025",
  "optionId": "OPT-FAST-1",
  "timeSlotId": "SLOT-101",
  "paymentStatus": "Paid",
  "shippingAmount": 3200000,
  "currency": "IRR"
}
\end{verbatim}
\end{latin}

\noindent نمونه پاسخ ثبت موفق سفارش لجستیک:
\begin{latin}
\begin{verbatim}
{
  "isSuccess": true,
  "logisticsOrderId": "LG-800025",
  "shipments": [
    {
      "shipmentId": "SHP-1001",
      "shipmentGroupId": "SG-1",
      "trackingCode": null,
      "status": "PendingAllocation"
    },
    {
      "shipmentId": "SHP-1002",
      "shipmentGroupId": "SG-2",
      "trackingCode": null,
      "status": "PendingAllocation"
    }
  ]
}
\end{verbatim}
\end{latin}

\begin{table}[htbp]
\centering
\begin{tabularx}{\textwidth}{l l X}
\toprule
\textbf{شناسه وضعیت} & \textbf{عنوان فارسی} & \textbf{توضیح} \\
\midrule
\lr{PendingAllocation} & در انتظار تخصیص & سفارش ثبت شده اما شرکت حمل یا راننده تخصیص نیافته است. \\
\lr{ReadyForPickup} & آماده برداشت & انبار کالا را بسته‌بندی کرده و آماده تحویل به پیک است. \\
\lr{PickedUp} & تحویل به حمل‌کننده & مرسوله توسط پیک از انبار تحویل گرفته شده است. \\
\lr{InTransit} & در مسیر & کالا در مسیر بین شهری یا انبار میانی است. \\
\lr{OutForDelivery} & در حال توزیع & مرسوله جهت تحویل نهایی به راننده توزیع واگذار شده است. \\
\lr{Delivered} & تحویل‌شده & سفارش با موفقیت به گیرنده تحویل داده شد. \\
\lr{DeliveryFailed} & ناموفق & تلاش برای تحویل ناموفق بوده و به انبار مرجوع می‌شود. \\
\lr{Canceled} & لغوشده & ارسال سفارش لغو گردیده است. \\
\bottomrule
\end{tabularx}
\caption{کدهای وضعیت مرسوله لجستیکی}
\end{table}

\section{خطاها، امنیت، قابلیت اطمینان و تکرارناپذیری}
در صورت بروز خطا، سیستم لجستیک اطلاعات دقیق خطا را در قالب پاسخ استاندارد ارسال خواهد کرد. جدول زیر نگاشت کدهای وضعیت \lr{HTTP} و کدهای خطای مربوطه را نشان می‌دهد:

\begin{table}[htbp]
\centering
\begin{tabularx}{\textwidth}{c l l X}
\toprule
\textbf{HTTP} & \textbf{کاربرد} & \textbf{کد خطا} & \textbf{توضیح} \\
\midrule
200 & موفق & \lr{QUOTE\_CREATED} & نقل‌قول با موفقیت ثبت شد. \\
\midrule
400 & نامعتبر & \lr{INVALID\_ADDRESS} & ساختار یا مقادیر آدرس ارسالی نامعتبر است. \\
\midrule
401/403 & احراز هویت & \lr{INVALID\_API\_KEY} & کلید دسترسی نامعتبر یا منقضی است. \\
\midrule
404 & عدم نگاشت & \lr{WAREHOUSE\_MAPPING\_NOT\_FOUND} & نگاشت کد انبار یا شهر یافت نشد. \\
\midrule
409 & منقضی شدن & \lr{QUOTE\_EXPIRED} & نقل‌قول منقضی گردیده است. \\
\midrule
422 & غیرقابل انجام & \lr{NO\_DELIVERY\_OPTION} & گزینه ارسالی معتبری وجود ندارد. \\
\midrule
429 & محدودیت & \lr{RATE\_LIMITED} & تعداد درخواست‌ها بیش از حد مجاز است. \\
\midrule
502/503 & اختلال سرویس & \lr{CARRIER\_UNAVAILABLE} & سیستم‌های حمل و نقل لجستیک موقتاً در دسترس نیستند. \\
\bottomrule
\end{tabularx}
\caption{کدهای وضعیت و خطای استاندارد سیستم لجستیک}
\end{table}

\paragraph{نمونه پاسخ خطای استاندارد (Standard Error Response):}
\begin{latin}
\begin{verbatim}
{
  "isSuccess": false,
  "statusCode": 422,
  "code": "NO_DELIVERY_OPTION",
  "message": "There is no active delivery option for the destination.",
  "traceId": "01J...",
  "details": [
    {
      "cartItemId": "502",
      "reasonCode": "OUT_OF_SERVICE_AREA",
      "message": "Bulky item is not deliverable to this city."
    }
  ]
}
\end{verbatim}
\end{latin}

\paragraph{الزامات فنی و امنیت ارتباطات:}
\begin{itemize}
    \item \textbf{خطاهای شبکه و موقت:} ردیابی درخواست‌ها با \lr{requestId} و \lr{traceId} انجام شده و مکانیزم تلاش مجدد (\lr{Retry}) صرفاً برای خطاهای موقت \lr{429}، \lr{502}، \lr{503} و \lr{Timeout} فعال می‌شود.
    \item \textbf{محدودیت زمانی استعلام (Timeout):} جهت تضمین تجربه کاربری روان، زمان انتظار استعلام گزینه‌ها حداکثر ۳ تا ۵ ثانیه خواهد بود و در صورت قطع شبکه، وضعیت‌های پیش‌فرض لوکال موقتاً فعال می‌شوند.
    \item \textbf{حفاظت از داده‌ها:} تمامی کلیدهای امنیتی مخفی نگاه داشته شده و اطلاعات حساس مشتریان در لاگ‌ها ماسک خواهند شد.
\end{itemize}

\section{سناریوهای نمونه کاربری}
\begin{itemize}
    \item \textbf{سناریو A (پیک اختصاصی تهران):} سبد شامل اقلام موجود در انبار مرکزی تهران و مقصد خریدار تهران است. وب‌سرویس لجستیک گزینه‌های پیک فوری، پیک زمان‌بندی‌شده و پست را برمی‌گرداند. کاربر بازه ۱۲ تا ۱۶ را انتخاب کرده و ظرفیت برای ۱۰ دقیقه رزرو می‌شود.
    \item \textbf{سناریو B (خارج از محدوده پیک):} آدرس خریدار به مشهد تغییر می‌کند. لجستیک پیک اختصاصی را حذف و صرفاً گزینه‌های پست و تیپاکس را با تاریخ تقریبی و هزینه‌های مربوطه نمایش می‌دهد.
    \item \textbf{سناریو C (سبد چندانبار مبدا):} کالای A در انبار تهران و کالای B در انبار کرج است. لجستیک سفارش را به دو مرسوله مستقل تفکیک کرده و هزینه کل ارسال را ۳۲۰,۰۰۰ ریال اعلام می‌کند. سفارش خریدار در قالب دو مرسوله مجزا تحویل خواهد شد.
    \item \textbf{سناریو D (تغییر تعداد اقلام):} خریدار در آخرین گام تعداد کالا را تغییر می‌دهد. کد هش سبد خرید (\lr{cartHash}) تغییر کرده، رزرو قبلی باطل شده و استعلام جدید فراخوانی می‌شود.
    \item \textbf{سناریو E (کالای حجیم غیرقابل ارسال):} کالا به علت حجم بالا صرفاً در محدوده تهران قابل تحویل است اما آدرس خریدار شیراز است. وب‌سرویس خطای \lr{422} با کد \lr{OUT\_OF\_SERVICE\_AREA} بازمی‌گرداند و سایت خریدار را برای حذف کالا یا تغییر آدرس راهنمایی می‌کند.
\end{itemize}

\section{سناریوهای تست و معیارهای پذیرش (UAT)}
تضمین صحت عملکرد سیستم لجستیک با اجرای سناریوهای تست زیر ارزیابی می‌گردد:
\begin{itemize}
    \item \textbf{T01:} تست استعلام تک‌کالایی از یک انبار به مقصد دارای پوشش پیک (نمایش بازه‌های زمانی تحویل).
    \item \textbf{T02:} تغییر مقصد به شهر فاقد پوشش پیک (مخفی شدن پیک و نمایش پست/تیپاکس).
    \item \textbf{T03:} استعلام سبد خرید چندانبار مبدا (بررسی صحت تفکیک مرسوله‌ها و جمع کل هزینه‌ها).
    \item \textbf{T04:} تست بروز خطای \lr{CITY\_MAPPING\_NOT\_FOUND} در صورت عدم نگاشت شهر مقصد در ادمین.
    \item \textbf{T05:} تست بروز خطای \lr{WAREHOUSE\_MAPPING\_NOT\_FOUND} در صورت عدم نگاشت کد انبار.
    \item \textbf{T06:} تست منقضی شدن رزرو نقل‌قول و بازگشت پیام \lr{QUOTE\_EXPIRED} در هنگام پرداخت.
    \item \textbf{T07:} تست پر شدن ظرفیت بازه زمانی و پیشنهاد بازه‌های جایگزین.
    \item \textbf{T08:} ابطال خودکار استعلام در صورت تغییر تعداد یا نوع کالاهای سبد خرید.
    \item \textbf{T09:} ارسال چندباره ثبت سفارش و اطمینان از عملکرد \lr{Idempotency-Key} جهت ایجاد تنها یک سفارش لجستیک.
    \item \textbf{T10:} قطع موقت ارتباط لجستیک و نمایش پیام تلاش مجدد همراه با ثبت شناسه پیگیری خطای \lr{TraceId}.
    \item \textbf{T11:} ارسال صحیح نام‌ها و کدهای خارجی استان‌ها و شهرها به صورت فارسی و استاندارد.
    \item \textbf{T12:} تفکیک فیلدهای \lr{warehouseId} و \lr{address} در پاسخ‌های دریافتی لجستیک.
\end{itemize}

\section{برنامه اصلاح ماژول فعلی و جداول دیتابیس پیشنهادی}
مراحل اجرای توسعه فنی به شرح زیر برنامه‌ریزی شده است:
\begin{enumerate}
    \item ممیزی کدهای افزونه ارسال فعلی و استخراج هدرها، جداول و هوک‌های تسویه حساب.
    \item اصلاح مدل‌های استان، شهر و انبار در ناپ‌کامرس و تعریف فیلد \lr{ExternalCode}.
    \item ایجاد جداول پیکربندی و نگاشت‌ها در دیتابیس ناپ‌کامرس.
    \item پیاده‌سازی کلاینت استعلام لجستیک (\lr{DeliveryQuoteClient}) همراه با منطق‌های \lr{Timeout}، \lr{Retry} و \lr{Logging}.
    \item ادغام با فرآیند ثبت و تسویه حساب سفارش و فراخوانی خودکار استعلام مجدد در صورت تغییر سبد یا آدرس.
    \item طراحی بخش نمایش گزینه‌ها و بازه‌های زمانی در فرانت‌سایت و ذخیره‌سازی انتخاب‌ها.
    \item پیاده‌سازی متد اعتبارسنجی نهایی پیش از ارجاع خریدار به درگاه پرداخت.
    \item ارسال خودکار و تکرارناپذیر اطلاعات نهایی سفارش پرداخت‌شده به سرور لجستیک.
    \item ایجاد صفحات مدیریت خطاها، تلاش مجدد دستی سفارشات لجستیک معلق و لاگ‌های خطا.
    \item اجرای آزمون‌های پذیرش کاربری (\lr{UAT}) در محیط استیج با شبیه‌سازی انواع سبدهای خرید.
\end{enumerate}

\begin{table}[htbp]
\centering
\begin{tabularx}{\textwidth}{l X}
\toprule
\textbf{نام جدول ماژول} & \textbf{کاربرد فنی و دیتابیس} \\
\midrule
\lr{LogisticsIntegrationSettings} & ذخیره‌سازی آدرس پایه وب‌سرویس، توکن دسترسی، زمان انتظار (\lr{Timeout}) و الگوهای تلاش مجدد. \\
\midrule
\lr{LogisticsLocationMapping} & نگاشت شناسه کشور، استان و شهرهای ناپ‌کامرس به کدهای سیستم لجستیک. \\
\midrule
\lr{LogisticsWarehouseMapping} & نگاشت شناسه‌های انبار ناپ‌کامرس به کدهای سیستمی انبار در لجستیک. \\
\midrule
\lr{LogisticsCartSelection} & ذخیره‌سازی مشخصات نقل‌قول فعال خریدار و بازه زمانی انتخاب‌شده کاربر. \\
\midrule
\lr{LogisticsOrderLink} & پیوند شناسه سفارش فروشگاه با شناسه سفارش در سیستم لجستیک و کدهای رهگیری مرسوله‌ها. \\
\midrule
\lr{LogisticsIntegrationLog} & لاگ خلاصه تراکنش‌ها، خطاها، پاسخ‌ها و شناسه‌های پیگیری (\lr{TraceId}) بدون نگهداری داده‌های حساس خریداران. \\
\bottomrule
\end{tabularx}
\caption{جداول دیتابیس ماژول لجستیک ناپ‌کامرس}
\end{table}

\newpage
\section{برآورد مالی و زمان‌بندی پیاده‌سازی ماژول H}
هزینه پیاده‌سازی نیازمندی‌های یکپارچه‌سازی لجستیک بر اساس چهار بخش کلیدی فوق به شرح زیر برآورد می‌گردد:

\begin{table}[htbp]
\centering
\begin{tabularx}{\textwidth}{c X c r}
\toprule
\textbf{ردیف} & \textbf{بخش / قابلیت نیازمندی لجستیک} & \textbf{زمان توسعه پیشنهادی} & \textbf{هزینه پیشنهادی (تومان)} \\
\midrule
۱ & طراحی لایه نگاشت داده‌ها و پنل مدیریت نگاشت‌ها & ۳ الی ۵ روز & ۷,۰۰۰,۰۰۰ \\
۲ & اتصال وب‌سرویس استعلام قیمت، تفکیک انبارها و بازه‌ها & ۱ هفته & ۱۲,۰۰۰,۰۰۰ \\
۳ & مکانیزم اعتبارسنجی مجدد سبد خرید پیش از پرداخت & ۳ روز & ۵,۰۰۰,۰۰۰ \\
۴ & ثبت تکرارناپذیر سفارش لجستیک و وب‌هوک‌های وضعیت & ۴ الی ۶ روز & ۸,۰۰۰,۰۰۰ \\
\midrule
\multicolumn{2}{r}{\textbf{جمع کل برآورد هزینه توسعه لجستیک (ماژول H)}} & \textbf{۳ الی ۴ هفته} & \textbf{۳۲,۰۰۰,۰۰۰} \\
\bottomrule
\end{tabularx}
\caption{جدول برآورد مالی و زمان‌بندی پیاده‌سازی بخش لجستیک}
\end{table}

\end{document}
